\newcommand*{\version}{3.32}

\documentclass[12pt,
              article,
              oneside
              ]{memoir}

\usepackage{iftex}

\ifluatex
  \usepackage{fontspec}
\fi

\usepackage[
  % latin,
  % english,
  spanish,% <- principal
  es-minimal,
  es-noshorthands,
  es-tabla,
]{babel}

\ifluatex
% LuaLaTeX --------------------
\defaultfontfeatures{Ligatures=TeX, Scale=MatchLowercase}
\usepackage[MainFeatures={medium,
  Ligatures=Common,
  Numbers=OldStyle,
  SmallCapsFeatures = {LetterSpace=5.0}}]
{junicodevf}
%
% \setsansfont{Libertinus Sans}
% (without the ligatures of TeX, so that -- stays two hyphens)
\setmonofont[Scale=MatchLowercase, Ligatures=TeXOff]{DejaVu Sans Mono}
% \setmonofont[Scale=MatchLowercase]{Libertinus Mono}
% \setmonofont[Scale=MatchLowercase}]{Fira Code}
\else
% PDFLaTeX ---------------------
\usepackage[T1]{fontenc}
\usepackage[osf]{libertinus}
\usepackage[scaled=0.9]{FiraMono}
% \usepackage[scaled=0.9]{cascadia-code}
% \usepackage[scaled=0.9]{roboto}
\input glyphtounicode
\pdfgentounicode=1
\fi

\usepackage[babel=true, verbose=false, tracking=true, expansion=true, protrusion=true, final, draft=false]{microtype}
\SetTracking{ encoding = *, shape = sc }{ 30 }
\SetTracking[context = notracking, ]{encoding = *}{0}

\usepackage[autostyle=false, style=british]{csquotes}
\usepackage{latexcolors}

\newcommand{\Rojo}{\color[HTML]{8B0000}}
\newcommand*{\rojo}[1]{\textcolor[HTML]{8B0000}{#1}}
\newcommand*{\rojoit}[1]{\textit{\textcolor[HTML]{8B0000}{#1}}}
\newcommand{\Azul}{\color{bluenivaca}}
\newcommand*{\azul}[1]{\textcolor{bluenivaca}{#1}}
\newcommand{\Verde}{\color{darkgreen}}
\newcommand*{\verde}[1]{\textcolor{darkgreen}{#1}}
\newcommand{\Amarillo}{\color{goldenpoppy}}
\newcommand*{\amarillo}[1]{\textcolor{goldenpoppy}{#1}}
\newcommand{\Morado}{\color{byzantium}}
\newcommand*{\morado}[1]{\textcolor{byzantium}{#1}}
\newcommand{\Naranja}{\color{orange}}
\newcommand*{\naranja}[1]{\textcolor{orange}{#1}}

% \usepackage[margin=3cm]{geometry}

\usepackage{graphicx}
\usepackage{booktabs}

\clubpenalty=9999
\widowpenalty=9999
\brokenpenalty=5000
\frenchspacing
\raggedbottom
\midsloppy
\parindent=14pt
\fboxsep=2pt

\renewcommand\thesection{\arabic{section}}
\chapterstyle{reparticle}

\usepackage{enumitem}
\setlist{noitemsep, nosep}
\setlist[description]{font=\itshape}

\newcommand*{\elipsis}{\,.\,.\,.\,}
\newcommand*{\elipsisb}{\textup{[\elipsis]}}

\renewenvironment{quote}%
 {\begin{list}{}%
              {\setlength\rightmargin{0pt}%
              \setlength\leftmargin{\parindent}}%
  \item[]\small\ignorespaces}
 {\unskip\end{list}}


\definecolor{Link}{HTML}{408080}
\usepackage{hyperref}
\hypersetup{
  colorlinks=true,
  linkcolor=blue,
  filecolor=blue,
  citecolor = black,
  urlcolor=Link,
}
\urlstyle{tt}
\usepackage{bookmark}

\setlength{\droptitle}{-2cm}

\title{\bfseries%
  \texttt{py-words} v.\version{} \\
  Manual de usuario}

\author{%
  Nicolás Vaughan\\
  \begin{minipage}[t]{0.8\linewidth}
    \centering\small%
    \href{mailto:nicolasvaughan@proton.me}{\texttt{nicolasvaughan@proton.me}}\\
    \url{https://nicolasvaughan.org}\\
    \url{https://github.com/nivaca/py-words}\\
    \url{https://codeberg.org/nivaca/py-words}\\
  \end{minipage}}


\date{2026-11-26}

% The pictures, made by tools/screenshots.py
\graphicspath{{img/}}
\usepackage{longtable}

% the title and the author in the PDF's own metadata
\hypersetup{pdftitle={py-words: manual de usuario}, pdfauthor={Nicolás Vaughan}}

% the index (makeindex, with the style manual.ist: see latexmkrc)
\makeindex

% menus, keys, and the pictures of the window
\newcommand*{\menu}[1]{\textit{#1}}
\newcommand*{\then}{\,>\,}
\newcommand*{\key}[1]{\texttt{#1}}
\newcommand*{\shot}[3]{%
  \begin{figure}[htbp]
    \centering
    \includegraphics[width=#2\textwidth,height=0.75\textheight,keepaspectratio]{#1}
    \caption{#3}
    \label{fig:#1}
  \end{figure}}
% the tilde of the interactive commands (~E, ~L ...)
\newcommand*{\tl}{\textasciitilde}

\begin{document}
\maketitle

\tableofcontents*

\section{Introducción}
\noindent
\texttt{py-words} es un diccionario latino que sabe gramática latina.
Analiza una palabra latina flexionada, o una línea entera de latín, en las
entradas del diccionario y las formas que puede ser, con sus significados
en español o en inglés, y encuentra el latín de una palabra española o
inglesa. Se puede usar de cuatro maneras: en una ventana, en la línea de
órdenes, en una terminal a pantalla completa y en un navegador web.
También está en línea, sin nada que instalar, en
\url{https://nivaca.pythonanywhere.com/}.

\index{WORDS (Whitaker)}\index{Lewis y Short}\index{De Miguel}\index{Echave-Sustaeta}\index{Latham}\index{Collatinus}%
Lo escribí en 2026. Su punto de partida fue WORDS, el programa de William
Whitaker, pero el código y los datos se rediseñaron, y el diccionario
creció: ahora tiene las palabras raras y las definiciones inglesas del
\textit{A Latin Dictionary} de Lewis y Short, las palabras, los nombres
propios y las definiciones españolas del \textit{Nuevo diccionario
latino-español etimológico} de R.~de~Miguel, el vocabulario básico de
J.~Echave-Sustaeta, la lista de palabras medievales de R.~E. Latham y las
vocales largas de Collatinus. Cada entrada dice de dónde viene
(sección~\ref{sec:codes}).

\index{convenciones}\index{macos y windows@macOS y Windows}%
Este manual explica primero cómo instalar y abrir el programa; luego, la
ventana, cómo buscar y cómo leer los resultados, y después los otros tres
programas, los ajustes y los editores del diccionario. Los menús van en
cursiva, con sus opciones tras un \enquote{>}, así:
\menu{Diccionario\then Ajustes\ldots}; las teclas, como \key{Ctrl+F}. En
macOS, \key{Ctrl} es la tecla Comando. Las imágenes muestran los programas
en español; con la interfaz en inglés, todos los menús, mensajes y textos
de ayuda están en inglés (sección~\ref{sec:languages}).

\index{licencia}%
\texttt{py-words} es software libre, con la Licencia Pública General de
GNU, versión 3 o posterior. Solo usa Python y su biblioteca estándar (con
Tk para las ventanas y curses para la terminal), así que no necesita nada
más; la lista completa del software que usa, con sus licencias, está en
\texttt{docs/LICENCIAS\_es.md}. Whitaker cedió WORDS y sus datos
\enquote{para cualquier uso}, y el texto de Lewis y Short viene de la
Perseus Digital Library, con la licencia CC BY-SA 4.0.

\section{Instalación}
\label{sec:install}

Cada versión se publica en GitHub y en Codeberg (las direcciones están en
la primera página), con estos archivos:

\index{appimage@AppImage}\index{fedora (copr, rpm)@Fedora (Copr, RPM)}\index{debian y ubuntu (deb)@Debian y Ubuntu (.deb)}\index{zip}%
\begin{itemize}
\item \texttt{py-words-v.VERSIÓN-windows-x64-setup.exe},
  \texttt{py-words-v.VERSIÓN-macos-arm64.pkg} (Apple silicon),
  \texttt{py-words-v.VERSIÓN-macos-x86\_64.pkg} (Intel) y
  \texttt{py-words-v.VERSIÓN-linux-x86\_64.run}: los instaladores, para
  Windows, macOS y Linux;
\item \texttt{py-words-v.VERSIÓN-x86\_64.AppImage}: la AppImage, para
  cualquier Linux;
\item \texttt{py-words-VERSIÓN-1.fc45.noarch.rpm}: el paquete de Fedora,
  firmado, con la clave que lo firma, \texttt{RPM-GPG-KEY-py-words};
\item \texttt{py-words\_VERSIÓN-1\_all.deb}: el paquete para Debian y
  Ubuntu, con su firma, \texttt{py-words\_VERSIÓN-1\_all.deb.asc};
\item \texttt{py-words-v.VERSIÓN.zip}: el programa mismo, para Linux,
  macOS y Windows.
\item \texttt{SHA256SUMS}, con \texttt{SHA256SUMS.asc}: las sumas SHA-256 de
  los archivos, y su firma. Cada archivo tiene además su propia firma al
  lado (\texttt{ARCHIVO.asc}), hecha con la misma clave, como la etiqueta de
  la versión; con la clave importada (\texttt{gpg --import
  RPM-GPG-KEY-py-words}), \texttt{gpg --verify ARCHIVO.asc ARCHIVO} comprueba
  un archivo.
\end{itemize}

\subsection{Qué se necesita}

\index{python@Python}\index{tk@Tk}%
Solo Python 3.13 o posterior. Las ventanas (el diccionario y su editor)
necesitan Tk, que viene con Python en Windows y con los instaladores de
python.org en macOS. En Linux, y con el Python de Homebrew en macOS, puede
ser un paquete aparte:

\begin{verbatim}
$ sudo dnf install python3-tkinter      (Fedora)
$ sudo apt install python3-tk           (Debian, Ubuntu)
$ brew install python-tk                (macOS, Homebrew)
\end{verbatim}

\noindent
\texttt{python3 --version} dice qué Python hay; si es demasiado antiguo,
el programa lo dice y se detiene. Las órdenes de este manual usan
\texttt{python3}; en Windows se escribe \texttt{python}. Los
instaladores, la AppImage y los paquetes traen o usan su propio Python,
así que no necesitan nada de esto.

\subsection{Los instaladores (Windows, macOS, Linux)}
\label{sec:installers}

\index{instaladores}\index{macos y windows@macOS y Windows}%
La manera más sencilla. Cada instalador trae su propio Python 3.13, con
Tk, e instala el programa, las entradas del menú, la orden
\texttt{py-words} (con el mismo primer argumento que la AppImage, más
abajo) y una manera de quitarlo todo. Una versión más nueva reemplaza a
la anterior; los ajustes, el historial y la caché se conservan. Como en
los paquetes, el diccionario instalado no se puede editar (la
sección~\ref{sec:editing} dice cómo hacerlo de todos modos).

\begin{description}
\item[Windows] (10 o posterior, 64 bits). Ejecute el \texttt{.exe}.
  Instala solo para usted, sin permisos de administrador, a menos que en
  su primera página elija instalar para todos los usuarios. Habla inglés
  o español. Añade la ventana al menú Inicio y, si quiere, la línea de
  órdenes, la versión de pantalla completa en la terminal y la página
  web, un acceso directo en el escritorio y la orden \texttt{py-words} en
  el PATH. Se quita desde «Aplicaciones instaladas» en la Configuración,
  y entonces pregunta si quitar también sus ajustes. No está firmado, así
  que SmartScreen puede advertir de él: «Más información» y luego
  «Ejecutar de todas formas».
\item[macOS] (11 o posterior). Abra el \texttt{.pkg} de su Mac (arm64
  para Apple silicon, x86\_64 para Intel). Pone \texttt{py-words.app} en
  Aplicaciones y, a menos que las deje fuera en «Personalizar»,
  \texttt{py-words (CLI)}, \texttt{(TUI)} y \texttt{(Web)}, que se abren
  en Terminal, y la orden \texttt{py-words}. No está firmado por Apple,
  así que macOS lo rechaza la primera vez: en Configuración del Sistema,
  «Privacidad y seguridad», elija «Abrir de todas formas».
  \texttt{sudo py-words uninstall} lo quita (\texttt{--purge}: también
  sus ajustes, su historial y la caché).
\item[Linux] (x86-64). Ejecute \texttt{sh py-words-v.VERSIÓN-linux-x86\_64.run}.
  Como usuario normal instala solo para usted, en
  \texttt{\textasciitilde/.local/share/py-words}, con la orden en
  \texttt{\textasciitilde/.local/bin}; con \texttt{sudo}, para todos, en
  \texttt{/opt/py-words} y \texttt{/usr/local/bin}. Añade las cuatro
  entradas del menú y los iconos. Primero dice qué va a hacer y
  pregunta: \texttt{--yes} no pregunta nada, \texttt{--only-gui} añade
  solo la entrada de la ventana, \texttt{--help} muestra las demás
  opciones. \texttt{py-words uninstall} lo quita.
\end{description}

\subsection{El ZIP, o una copia del repositorio}

\index{zip}\index{git}%
Se descomprime el ZIP en cualquier parte, o se clona el repositorio:

\begin{verbatim}
$ git clone https://github.com/nivaca/py-words.git
\end{verbatim}

\noindent
La carpeta es el programa entero: no hay nada que compilar ni instalar, y
funciona desde donde está.

\begin{verbatim}
$ cd py-words
$ python3 words.py amicus          (una primera prueba)
$ python3 words_gui.py             (la ventana)
\end{verbatim}

\subsection{La AppImage (Linux)}

\index{appimage@AppImage}%
La AppImage es un solo archivo que trae su propio Python, con Tk. Se hace
ejecutable y se abre:

\begin{verbatim}
$ chmod +x py-words-v.3.32-x86_64.AppImage
$ ./py-words-v.3.32-x86_64.AppImage              (la ventana)
$ ./py-words-v.3.32-x86_64.AppImage cli amicus   (en la terminal)
$ ./py-words-v.3.32-x86_64.AppImage help         (los demás)
\end{verbatim}

\noindent\index{editores del diccionario!datos de solo lectura}%
El primer argumento elige el programa: \texttt{gui} (el de omisión),
\texttt{cli}, \texttt{tui}, \texttt{web}, \texttt{edit},
\texttt{edit-gui}, \texttt{check} o \texttt{convert}; lo que sigue va a
ese programa. En la AppImage no se puede escribir, así que en ella no se
puede editar el diccionario (la sección~\ref{sec:editing} explica cómo
hacerlo de todos modos). Herramientas como AppImageLauncher o Gear Lever
la añaden al menú del sistema.

\subsection{Fedora}

\index{fedora (copr, rpm)@Fedora (Copr, RPM)}%
Lo más sencillo es el repositorio de Copr. El paquete usa el Python y el
Tk del propio Fedora, y añade cuatro entradas al menú del sistema, con el
icono:

\begin{verbatim}
$ sudo dnf copr enable nivaca/py-words
$ sudo dnf install py-words
\end{verbatim}

\noindent
Las actualizaciones llegan luego con \texttt{sudo dnf upgrade}. El RPM de
cada versión también se puede instalar a mano, tras importar una vez su
clave con \texttt{sudo rpm --import RPM-GPG-KEY-py-words}. Funciona en
cualquier Fedora con Python 3.13 o posterior, sea cual sea el número
\texttt{fc} de su nombre. La orden \texttt{py-words} abre entonces los
programas, con el mismo primer argumento que la AppImage:

\begin{verbatim}
$ py-words                 (la ventana)
$ py-words cli amicus      (la línea de órdenes)
$ py-words help            (los demás programas)
\end{verbatim}

\noindent\index{editores del diccionario!datos de solo lectura}%
El programa se instala en \texttt{/usr/share/py-words}, donde no se puede
escribir, así que tampoco ahí se puede editar el diccionario, como en la
AppImage.

\subsection{Debian y Ubuntu}

\index{debian y ubuntu (deb)@Debian y Ubuntu (.deb)}%
El paquete \texttt{.deb} instala lo mismo que el RPM, y también usa el
Python y el Tk del sistema. Necesita Python 3.13 o posterior (Debian 13,
Ubuntu 25.04 o posteriores); \texttt{apt} instala con él
\texttt{python3-tk}:

\begin{verbatim}
$ sudo apt install ./py-words_3.32-1_all.deb
\end{verbatim}

\noindent
Su firma viene aparte, en el archivo \texttt{.deb.asc}, hecha con la misma
clave que el RPM; \texttt{gpg --verify} la comprueba una vez importada la
clave. La orden \texttt{py-words} abre entonces los programas, como con el
RPM, y \texttt{sudo apt remove py-words} lo quita. No hay repositorio:
cada versión nueva se instala a mano de la misma manera.

\subsection{Una entrada en el menú}

\index{entradas del menu@entradas del menú}\index{macos y windows@macOS y Windows}%
Con el ZIP o una copia del repositorio, basta ejecutar una vez

\begin{verbatim}
$ python3 words_install.py
\end{verbatim}

\noindent
para añadir el programa al menú del sistema, con cuatro entradas: la
ventana, la línea de órdenes y la versión de terminal a pantalla completa
(estas dos se abren en una terminal), y la versión web (su servidor corre
en una terminal, y se abre el navegador). En Linux van al menú de
aplicaciones; en Windows, al menú Inicio, y en macOS, a
\texttt{\textasciitilde/Applications}. \texttt{--only-gui} añade solo la
ventana, y \texttt{--uninstall} las quita. No necesita permisos de
administrador. Las entradas abren el programa en su carpeta, así que si la
carpeta se mueve hay que ejecutarlo de nuevo.

\subsection{Actualizar y desinstalar}

\index{actualizar}\index{desinstalar}%
Con git, \texttt{git pull} en la carpeta; si no, se descarga el ZIP de
nuevo y se reemplaza la carpeta. Si se editó el diccionario, conviene
guardar antes una copia de \texttt{data/dictionary.xml}. Para
desinstalarlo, se quitan las entradas del menú con \texttt{python3
words\_install.py --uninstall} y se borra la carpeta; los ajustes, el
historial y la caché también pueden borrarse (sección~\ref{sec:files}).

\section{Los cuatro programas}
\label{sec:programs}

\index{programas}\index{ventana}\index{linea de ordenes@línea de órdenes}\index{terminal a pantalla completa}\index{pagina web@página web}%
El mismo diccionario tiene cuatro caras, que comparten los ajustes y el
historial de búsquedas:

\begin{longtable}{@{}lp{0.62\textwidth}@{}}
\toprule
Programa & Qué es\\
\midrule
\endhead
\texttt{words\_gui.py} & la ventana (sección~\ref{sec:window})\\
\texttt{words.py} & la línea de órdenes, con un modo interactivo (sección~\ref{sec:cli})\\
\texttt{words\_tui.py} & la ventana, pero en una terminal, a pantalla completa (sección~\ref{sec:tui})\\
\texttt{words\_web.py} & un pequeño servidor web, y el diccionario en el navegador (sección~\ref{sec:web})\\
\texttt{meanings.py} & \texttt{words.py} solo con las formas de diccionario y los significados\\
\texttt{words\_edit\_gui.py} & el editor del diccionario, en una ventana (sección~\ref{sec:editing})\\
\texttt{words\_edit.py} & el editor del diccionario, en la terminal\\
\texttt{words\_check.py} & revisa los archivos de datos tras editarlos a mano\\
\texttt{words\_install.py} & añade el programa al menú del sistema\\
\bottomrule
\end{longtable}

\noindent
Desde la carpeta del programa se abren con \texttt{python3}:

\begin{verbatim}
$ python3 words_gui.py                     (una ventana)
$ python3 words_gui.py perspicio           (que la busca)
$ python3 words_gui.py -e mirar            (una ventana, al latín)
$ python3 words.py rem acu tetigisti       (buscar una línea)
$ python3 words.py                         (interactivo)
$ python3 words_tui.py                     (en la terminal)
$ python3 words_web.py                     (en el navegador)
\end{verbatim}

\noindent\index{version en linea@versión en línea}%
La versión en línea, \url{https://nivaca.pythonanywhere.com/}, es la
página de \texttt{words\_web.py} en un servidor público. Allí cada
visitante tiene sus propios ajustes y búsquedas anteriores, guardados en
las cookies de su navegador, y nada de eso se guarda en el servidor.

\section{La ventana}
\label{sec:window}

\shot{gui_es}{1}{La ventana, en modo oscuro, tras buscar \textit{perspicio}.}

\index{ventana}\index{cuadro de busqueda@cuadro de búsqueda}\index{buscar}%
La figura~\ref{fig:gui_es} muestra la ventana. Arriba está el cuadro donde
se escriben las palabras, y \menu{Buscar}, que las busca (\key{Intro} hace
lo mismo). El cuadro guarda las búsquedas anteriores: \key{Abajo}, o su
flecha, las lista, la más reciente primero, y elegir una la busca de
nuevo.

\index{buscar en}\index{direccion@dirección}\index{significados en}\index{busqueda aproximada@búsqueda aproximada}%
Debajo están las opciones de la búsqueda:

\begin{description}
\item[Buscar en] \menu{Local}, este diccionario, o \menu{Logeion}, los
  diccionarios de la Universidad de Chicago (sección~\ref{sec:logeion}),
  con la lista de los diccionarios de Logeion al lado;
\item[la dirección] \menu{Latín\,$\rightarrow$\,español} o
  \menu{español\,$\rightarrow$\,latín} (con los significados en inglés,
  \menu{inglés} en lugar de \menu{español});
\item[Significados en] el idioma de los significados
  (sección~\ref{sec:languages});
\item[Búsqueda aproximada] cuando una palabra no encuentra nada, los
  resultados de sus grafías cercanas (sección~\ref{sec:fuzzy}).
\end{description}

\noindent
Estas opciones se guardan: la próxima vez, cualquiera de los programas
empieza como se dejó.

\index{resultados!enlaces}\index{resultados!explicaciones}\index{doble clic}\index{clic derecho}%
Los resultados ocupan el resto de la ventana. Son el mismo texto que en la
línea de órdenes, dispuesto para la vista: la palabra buscada va más
grande, la forma de diccionario en negrita con su lema en color, los
códigos y las notas en gris, y los análisis de las formas en azul. La
primera palabra de una entrada es un enlace que la busca, y un doble clic
en cualquier palabra también la busca (una palabra de los significados se
busca en latín, en la otra dirección). Con el puntero sobre los códigos
entre corchetes o sobre un análisis (\texttt{GEN.PL}), un recuadro explica
qué significan. El clic derecho copia o busca una palabra.

\index{linea de estado@línea de estado}\index{modo oscuro}%
La línea de abajo dice qué está haciendo el programa. A su derecha, la
luna oscurece las ventanas (el diccionario, sus diálogos y el editor), y
luego el sol las vuelve a aclarar. A su lado, el código del idioma de la
interfaz (EN, ES o LA) abre la lista de los idiomas.

\subsection{Los menús}

\index{menus@menús}%
\begin{description}
\item[Diccionario] \menu{Editar diccionario\ldots} abre el editor
  (sección~\ref{sec:editing}); \menu{Ajustes\ldots} muestra todos los
  ajustes (sección~\ref{sec:settings}); \menu{Salir}.
\item[Editar] \menu{Cortar}, \menu{Copiar} y \menu{Pegar}, en el cuadro de
  búsqueda o en los resultados.
\item[Ver] \menu{Buscar en los resultados\ldots}, el tamaño del texto, las
  fuentes y el \menu{Modo oscuro}.
\item[Ayuda] \menu{Ayuda}, la guía para buscar y para leer los
  resultados; \menu{Sitio web}, la página del programa; \menu{Acerca de},
  la versión y los créditos.
\end{description}

\subsection{Buscar un texto en los resultados}

\index{buscar en los resultados}%
\menu{Ver\then Buscar en los resultados\ldots} (\key{Ctrl+F}) abre una
pequeña barra bajo los resultados. Lo que se escribe allí se resalta a
medida que se escribe; \key{Intro} va al lugar siguiente,
\key{Mayús+Intro} al anterior, y \key{Esc} cierra la barra.

\subsection{Texto y fuentes}

\index{fuentes}\index{tamano del texto@tamaño del texto}\index{pantallas de alta resolucion@pantallas de alta resolución}%
\key{Ctrl++} y \key{Ctrl+-} agrandan y achican el texto, y \key{Ctrl+0}
devuelve el tamaño habitual. \menu{Ver\then Fuente\ldots} elige las
fuentes: una de ancho fijo para los resultados, porque sus columnas deben
quedar alineadas, y otra para el resto de la ventana. Ambas se guardan.
En una pantalla de alta resolución, las ventanas siguen el escalado del
escritorio.

\section{Buscar}
\label{sec:searching}

\subsection{Desde el latín}

\index{buscar!desde el latin@desde el latín}\index{mayusculas@mayúsculas}\index{u y v, i y j}\index{marcas de cantidad}%
Se escribe una palabra latina, o una línea entera de latín: cada palabra
se analiza por turno. No importan las mayúsculas, ni la u y la v, ni la i
y la j (\textit{iulius}, \textit{Julius} e \textit{IVLIVS} son la misma
palabra), ni las marcas de cantidad y otros acentos: \textit{amīcus} es
\textit{amicus}, y \textit{cælum} es \textit{caelum}.

\index{trucos}\index{prefijos y sufijos}\index{encliticos@enclíticos}\index{grafias medievales@grafías medievales}%
Una palabra se busca primero como una raíz y una desinencia del
diccionario. Si así no se encuentra nada, el programa prueba prefijos,
sufijos y enclíticos (\textit{-que}, \textit{-ne}, \textit{-ve}), y luego
otras grafías: un perfecto sincopado, dos palabras escritas como una
(\textit{rempublicam}), una \textit{-st} final por \textit{est}
(\textit{factumst}) y las grafías del latín medieval (\textit{e} por
\textit{ae}, \textit{y} por \textit{i}, \textit{michi} por
\textit{mihi}). Siempre dice lo que hizo:

\begin{verbatim}
hystorie
  ortografía: Grafía medieval hystorie -> historiae
    'hystorie' puede ser una grafía medieval de historiae
  historia, historiae
      sustantivo, 1.ª declinación, femenino  [XXXAX]
    Historia; Historieta, fábula, cuento; (De Miguel)
    histori-ae     GEN.SG
\end{verbatim}

\noindent\index{nombres propios}%
Una palabra desconocida con mayúscula inicial se toma como un nombre, y se
deja como está.

\subsection{Hacia el latín}

\index{buscar!hacia el latin@hacia el latín}\index{categorias gramaticales@categorías gramaticales}%
En la otra dirección (\menu{español\,$\rightarrow$\,latín} en la ventana,
\texttt{-e} en la línea de órdenes), se escribe una palabra del idioma de
los significados: se listan las entradas latinas cuyos significados la
tienen, las mejores primero, seis como mucho, y luego cuántas más hay. Una
categoría gramatical tras la palabra deja solo esa: \texttt{amigo ADJ}.
Las categorías son \texttt{N}, \texttt{PRON}, \texttt{ADJ}, \texttt{NUM},
\texttt{ADV}, \texttt{V}, \texttt{PREP}, \texttt{CONJ} e \texttt{INTERJ}.
Las palabras muy comunes (\textit{el}, \textit{la}, \textit{de}) nunca se
buscan, y los acentos no importan: \textit{corazon} encuentra
\textit{corazón}.

\subsection{La búsqueda aproximada}
\label{sec:fuzzy}

\index{busqueda aproximada@búsqueda aproximada}\index{erratas}%
Por lo general, una palabra con una errata no encuentra nada. Con la
búsqueda aproximada, a una palabra de tres letras o más que sigue sin
encontrar nada le siguen los resultados de las palabras a un cambio de
ella (una letra omitida, añadida o cambiada, o dos letras intercambiadas),
tres como mucho, primero las de las entradas más frecuentes:

\begin{verbatim}
amicuss
  (desconocida)
  (grafías cercanas: amicus)

amicus
  amicus, amici  sustantivo, 2.ª declinación, masculino  [XXXBO]
  ...
\end{verbatim}

\noindent
La búsqueda hacia el latín hace lo mismo con las palabras de los
significados: \textit{corazn} encuentra \textit{corazón}, y
\textit{amgio}, \textit{amigo}. Está activada mientras no se desactive:
\menu{Búsqueda aproximada} en la ventana y en la página web, \key{F3} en
la terminal, o \texttt{--no-fuzzy} en la línea de órdenes, para una sola
ejecución.

\section{Leer los resultados}
\label{sec:results}

\index{resultados}%
Esto es lo que \texttt{py-words} muestra para \textit{amicorum}:

\begin{verbatim}
amicorum
  amicus, amici  sustantivo, 2.ª declinación, masculino  [XXXBO]
    El amigo (en el sent. más general de la palabra); ...
    amīc-ōrum      GEN.PL
  amicus, amica -um, amicior -or -us, amicissimus -a -um
      adjetivo, 1.ª/2.ª declinación  [XXXAO]
    Amigo, amigable, amistoso, favorable, ... (De Miguel)
    amīc-ōrum      GEN.PL.M
    amīc-ōrum      GEN.PL.N
\end{verbatim}

\index{resultados!lineas de@líneas de}\index{partes principales}\index{vocales largas}%
\begin{itemize}
\item La primera línea es la palabra buscada. Detrás, con sangría, cada
  entrada que puede ser.
\item La primera línea de una entrada tiene su forma de diccionario (las
  partes principales: \textit{amicus, amici}), su categoría y su gramática
  en palabras (sustantivo, 2.ª declinación, masculino), y los códigos
  entre corchetes (sección~\ref{sec:codes}); algunas entradas añaden su
  época o su frecuencia en palabras (arcaico, menos frecuente\ldots).
\item La línea siguiente es el significado. Sus acepciones van separadas
  por punto y coma, la más común primero. Al final, el diccionario del que
  viene, cuando no es WORDS: (De Miguel), (Echave-Sustaeta), (Lewis and
  Short). \texttt{[en]} ante un significado indica que no tiene
  traducción, y se da en inglés.
\item Luego vienen las formas que puede ser la palabra: la raíz y la
  desinencia, unidas por un guion, con las vocales largas marcadas donde
  se conocen (\textit{amīc-ōrum}), y qué forma es, con las abreviaturas de
  la sección~\ref{sec:glosses}. Una misma grafía puede ser varias formas:
  el \textit{amīc-ōrum} del adjetivo es masculino y neutro.
\item Una línea antes de una entrada puede decir cómo se halló la palabra:
  un prefijo, un sufijo o un enclítico, con su significado, o un cambio de
  grafía.
\end{itemize}

\noindent\index{tiempos compuestos}%
Un participio seguido de una forma de \textit{esse} se lee como un tiempo
compuesto:

\begin{verbatim}
laudatus
  laudo, laudare, laudavi, laudatus
      verbo, 1.ª conjugación  [XXXAX]
    Alabar, bonos; ponderar la excelencia de (en medic.), ...
    laudāt-us      PRF.PASS.PTCP.NOM.SG.M
    participio + est  3SG.PRF.PASS.IND
      participio perfecto pasivo + verbo «ser»
        => sistema de perfecto pasivo
\end{verbatim}

\noindent\index{resultados!notas}%
Algunas notas van entre paréntesis:

\begin{longtable}{@{}>{\raggedright}p{0.4\textwidth}p{0.54\textwidth}@{}}
\toprule
Nota & Qué dice\\
\midrule
\endhead
(desconocida) & la palabra no se encontró\\
(grafías cercanas: \ldots) & siguen los resultados de la búsqueda aproximada\\
(no se muestran las formas menos probables) & el ajuste \texttt{trim\_output} omitió algunas\\
(sin resultados) & la búsqueda hacia el latín no encontró nada\\
(palabras cercanas: \ldots) & la búsqueda aproximada hacia el latín\\
(no se muestran \ldots\ más) & tras las seis mejores entradas, hacia el latín\\
\bottomrule
\end{longtable}

\subsection{Los códigos entre corchetes}
\label{sec:codes}

\index{codigos entre corchetes@códigos entre corchetes}\index{epoca@época}\index{area@área}\index{region@región}\index{frecuencia}\index{fuentes de las entradas}%
Las cinco letras entre corchetes que siguen a una entrada dicen cuándo,
dónde, con qué frecuencia y según qué autoridad se usa la palabra, en este
orden: su época, su área (el tema al que pertenece), su región, su
frecuencia y su fuente. Una X solo dice que no hay nada particular que
señalar. Así, el \texttt{[XXXBO]} de \textit{amicus} se lee: en uso en
todas las épocas, general, de todas partes, frecuente, y tomada del Oxford
Latin Dictionary. Una entrada puede tener más de una fuente: las demás
siguen a las cinco letras, cada una tras un \texttt{+}
(\texttt{[XXXBO+A]} viene también de Echave-Sustaeta). La época, si se
conoce, y la frecuencia, si es menor o más rara, también se escriben en
palabras tras los códigos. Las letras son estas:

\begin{longtable}{@{}lp{0.8\textwidth}@{}}
\toprule
Letra & Época\\
\midrule
\endhead
X & en uso en todas las épocas\\
A & arcaico\\
B & antiguo (antes del periodo clásico)\\
C & clásico\\
D & tardío (después del 200)\\
E & posterior (después del 600)\\
F & medieval (después del 1000)\\
G & erudito (Renacimiento y después)\\
H & moderno\\
\bottomrule
\end{longtable}

\begin{longtable}{@{}lp{0.8\textwidth}@{}}
\toprule
Letra & Área\\
\midrule
\endhead
X & general\\
A & agricultura, plantas, animales, la tierra\\
B & biología, medicina, el cuerpo\\
D & teatro, música, las artes\\
E & la Iglesia, la Biblia\\
G & gramática, retórica, lógica, literatura, escuela\\
L & derecho, gobierno, dinero, títulos\\
P & poesía\\
S & ciencia, filosofía, matemáticas, medidas\\
T & técnica, arquitectura, agrimensura\\
W & guerra, el ejército, la marina\\
Y & mitología\\
\bottomrule
\end{longtable}

\begin{longtable}{@{}lp{0.8\textwidth}@{}}
\toprule
Letra & Región\\
\midrule
\endhead
X & en todas partes\\
A & África\\
B & Britania\\
C & China\\
D & Escandinavia\\
E & Egipto\\
F & Francia, la Galia\\
G & Alemania\\
H & Grecia\\
I & Italia, Roma\\
J & India\\
K & los Balcanes\\
N & los Países Bajos\\
P & Persia\\
Q & el Oriente Próximo\\
R & Rusia\\
S & Hispania\\
U & Europa oriental\\
\bottomrule
\end{longtable}

\begin{longtable}{@{}lp{0.8\textwidth}@{}}
\toprule
Letra & Frecuencia\\
\midrule
\endhead
X & desconocida\\
A & muy frecuente\\
B & frecuente\\
C & común\\
D & menos frecuente\\
E & poco común\\
F & muy rara\\
I & solo en inscripciones\\
M & solo en grafitos\\
N & solo en Plinio\\
\bottomrule
\end{longtable}

\begin{longtable}{@{}lp{0.8\textwidth}@{}}
\toprule
Letra & Fuente\\
\midrule
\endhead
X & general o desconocida\\
A & Echave-Sustaeta, \textit{Lengua latina. Vocabulario básico}\\
B & Beeson, \textit{Medieval Latin}\\
C & \textit{Cassell's Latin Dictionary}\\
D & Adams, \textit{Latin Sexual Vocabulary}\\
E & Stelten, \textit{Ecclesiastical Latin}\\
F & Deferrari, santo Tomás de Aquino\\
G & Gildersleeve y Lodge, \textit{Latin Grammar}\\
H & Collatinus (Ouvrard)\\
I & Leverett, \textit{Lexicon}\\
J & Bracton, \textit{De Legibus}\\
K & \textit{Calepinus Novus} (Licoppe)\\
L & Lewis, \textit{Elementary Latin Dictionary}\\
M & Latham, \textit{Medieval Word List}; DMLBS\\
N & Nelson, lista de palabras\\
O & \textit{Oxford Latin Dictionary}\\
P & Souter, \textit{Later Latin}\\
Q & otros diccionarios\\
R & Plater y White, \textit{Vulgate grammar}\\
S & Lewis y Short, \textit{A Latin Dictionary}\\
T & hallada en una traducción\\
U & Du Cange\\
V & Blatt, Saxo\\
W & conjetura de Whitaker\\
Y & provisional\\
Z & enviada por un usuario\\
m & De Miguel, \textit{Nuevo diccionario latino-español etimológico}\\
\bottomrule
\end{longtable}

\subsection{Las abreviaturas de los análisis}
\label{sec:glosses}

\index{abreviaturas de glosa}\index{Reglas de Glosado de Leipzig}%
Las formas se analizan con las abreviaturas de las Reglas de Glosado de
Leipzig, unidas por puntos; la persona y el número van juntos, así que
\texttt{2SG.PRF.ACT.IND} es la 2.ª persona del singular del perfecto de
indicativo activo. Las marcadas con * no están en la lista de Leipzig.
Las abreviaturas son las mismas en los dos idiomas. Con el ajuste
\texttt{use\_glosses} desactivado (\texttt{--no-glosses} en la línea de
órdenes), los análisis se escriben en palabras: \enquote{acusativo
singular}.

\begin{longtable}{@{}lp{0.6\textwidth}@{}}
\toprule
Abreviatura & Significado\\
\midrule
\endhead
\texttt{NOM}, \texttt{GEN}, \texttt{DAT} & nominativo, genitivo, dativo\\
\texttt{ACC}, \texttt{ABL}, \texttt{VOC}, \texttt{LOC} & acusativo, ablativo, vocativo, locativo\\
\texttt{SG}, \texttt{PL} & singular, plural\\
\texttt{M}, \texttt{F}, \texttt{N}, \texttt{C}* & masculino, femenino, neutro, común\\
\texttt{1}, \texttt{2}, \texttt{3} & primera, segunda, tercera persona\\
\texttt{PRS}, \texttt{IMPF}*, \texttt{FUT} & presente, imperfecto, futuro\\
\texttt{PRF}, \texttt{PLUP}*, \texttt{FUT.PRF}* & perfecto, pluscuamperfecto, futuro perfecto\\
\texttt{ACT}*, \texttt{PASS} & activo, pasivo\\
\texttt{IND}, \texttt{SBJV}, \texttt{IMP} & indicativo, subjuntivo, imperativo\\
\texttt{INF}, \texttt{PTCP}, \texttt{SUP}* & infinitivo, participio, supino\\
\texttt{POS}*, \texttt{CMPR}*, \texttt{SPRL}* & positivo, comparativo, superlativo\\
\texttt{CARD}*, \texttt{ORD}* & numeral cardinal, ordinal\\
\texttt{DISTR}, \texttt{MULT}* & numeral distributivo, multiplicativo (adverbio numeral)\\
\bottomrule
\end{longtable}

\section{Los idiomas}
\label{sec:languages}

\subsection{El idioma de los significados}

\index{significados en}\index{significados en espanol@significados en español}\index{De Miguel}\index{Echave-Sustaeta}%
Los significados están en español o en inglés, y también el resto de los
resultados (\enquote{sustantivo, 2.ª declinación, masculino}). Los
significados en español vienen del diccionario de De Miguel siempre que
este tiene la palabra, tal como él los escribió en 1867, con su ortografía
(\textit{razon}, \textit{á}), y lo dicen al final, \enquote{(De Miguel)};
unos pocos vienen de Echave-Sustaeta; los demás se tradujeron del inglés.
Se elige con \menu{Significados en} en la ventana y en la página web, con
\key{F5} en la terminal, o con \texttt{--language en} en la línea de
órdenes, para una sola ejecución. La búsqueda hacia el latín parte del
idioma elegido.

\subsection{El idioma de la interfaz}

\index{idioma de la interfaz}%
Los programas hablan español, inglés o latín medieval: sus menús,
botones, mensajes,
ayuda y ajustes, los editores y el revisor. Se elige con el botón de su
código (EN, ES o LA), abajo a la derecha de la ventana y arriba en la
página web, con \menu{Idioma de la interfaz} en los ajustes de la ventana
y de la página web, con \key{F6} en la terminal (su código está a la
derecha de la línea de arriba), o al final del diálogo \texttt{\#} de la línea de
órdenes (\texttt{--interface en} o \texttt{--interface la} para una sola
ejecución). La ventana y la página web cambian en el acto. Cambiar la
interfaz cambia también los significados al mismo idioma, si el
diccionario lo tiene; no hay significados en latín, así que con la
interfaz latina siguen como estaban, y también las notas de los
resultados, como \enquote{(unknown)}; después, cada uno puede elegirse por
separado, de modo que la interfaz puede estar en español y los
significados en inglés.

\section{Logeion}
\label{sec:logeion}

\index{Logeion}%
En lugar de las entradas de este diccionario, los programas pueden
mostrar las de Logeion (\url{https://logeion.uchicago.edu}), los
diccionarios latinos y griegos de la Universidad de Chicago, con la
definición breve, la frecuencia y ejemplos de su corpus. Necesita
internet. Se elige con \menu{Buscar en} en la ventana y en la página web,
\key{F8} en la terminal, \texttt{\tl G} en la línea de órdenes
interactiva, o \texttt{--logeion} para una sola ejecución.

\index{lemas}%
Logeion no analiza las formas flexionadas. Por eso las palabras se
analizan aquí primero, y a Logeion se le piden sus lemas: \textit{dixit}
muestra \textit{dico}, y \textit{est} muestra \textit{edo} y \textit{sum}.
Cada palabra muestra primero su análisis, sin los significados, y luego
las páginas de Logeion de sus lemas. Mientras Logeion está elegido, la
búsqueda parte solo del latín, y la otra dirección, el idioma de los
significados, la búsqueda aproximada y el editor quedan en gris.

\index{Logeion!diccionarios}\index{DMLBS}%
Las entradas pueden ser las de un diccionario de Logeion, o las de todos;
mientras no se elija otro, es el DMLBS. Se elige con la lista que está
junto a \menu{Buscar en}, con \key{F5} en la terminal (mientras Logeion
está elegido), con \texttt{\tl D} en la línea de órdenes interactiva, o
con \texttt{--dictionary} y su código:

\begin{longtable}{@{}lp{0.7\textwidth}@{}}
\toprule
Código & Diccionario\\
\midrule
\endhead
\texttt{all} & todos los diccionarios\\
\texttt{ls} & Lewis \& Short\\
\texttt{lewis} & Lewis, \textit{Elementary Latin Dictionary}\\
\texttt{georges} & Georges\\
\texttt{gaffiot} & Gaffiot 2016\\
\texttt{dmlbs} & DMLBS (\textit{Dictionary of Medieval Latin from British Sources})\\
\texttt{lane} & LaNe\\
\texttt{ducange} & Du Cange\\
\texttt{lana} & Lana-Michelone\\
\texttt{frieze} & Frieze-Dennison (Virgilio)\\
\texttt{sinicum} & Latino-Sinicum\\
\bottomrule
\end{longtable}

\noindent
Al cambiar de diccionario, los resultados se muestran de nuevo en el
acto, sin volver a preguntar a Logeion. En la ventana, los enlaces de una
entrada a otras palabras buscan esas palabras, y los demás enlaces se
abren en el navegador.

\section{La línea de órdenes}
\label{sec:cli}

\shot{cli_es}{1}{La línea de órdenes, interactiva, con \textit{perspicio}.}

\index{linea de ordenes@línea de órdenes}%
\texttt{words.py} busca las palabras que se le dan, como una línea de
latín, o, sin palabras, abre el modo interactivo
(figura~\ref{fig:cli_es}). \texttt{meanings.py} es lo mismo que
\texttt{words.py -m}.

\begin{verbatim}
$ python3 words.py rem acu tetigisti
$ python3 words.py -e amigo --pos N
$ python3 words.py -f texto.txt -o resultados.txt
$ python3 words.py --language en amicorum
$ python3 words.py --logeion arma
$ python3 words.py --dictionary gaffiot arma
\end{verbatim}

\index{linea de ordenes@línea de órdenes!opciones}%
\begin{longtable}{@{}>{\raggedright}p{0.36\textwidth}p{0.58\textwidth}@{}}
\toprule
Opción & Qué hace\\
\midrule
\endhead
\texttt{-f ARCHIVO} & busca cada línea de ARCHIVO (\texttt{-} para la entrada estándar)\\
\texttt{-o ARCHIVO} & escribe los resultados en ARCHIVO\\
\texttt{-e} & del idioma de los significados al latín\\
\texttt{--pos CAT} & con \texttt{-e}, solo esta categoría gramatical\\
\texttt{-m} & solo las formas de diccionario y los significados\\
\texttt{--no-glosses} & los análisis en palabras, no en abreviaturas\\
\texttt{--no-fuzzy}, \texttt{--fuzzy} & sin o con la búsqueda aproximada\\
\texttt{--logeion}, \texttt{--local} & Logeion, o este diccionario\\
\texttt{--dictionary DIC} & el diccionario de Logeion (sección~\ref{sec:logeion})\\
\texttt{--dictionaries} & lista los diccionarios de Logeion\\
\texttt{--language CÓDIGO} & el idioma de los significados (\texttt{es}, \texttt{en})\\
\texttt{--languages} & lista los idiomas de los significados\\
\texttt{--interface CÓDIGO} & el idioma de la interfaz\\
\texttt{--datadir CARPETA} & los archivos de datos de CARPETA\\
\texttt{--guide} & imprime la ayuda: cómo leer los resultados\\
\texttt{--config} & crea el archivo de configuración y dice dónde está\\
\texttt{--clear-history} & olvida las búsquedas anteriores\\
\texttt{--version} & la versión y los créditos\\
\bottomrule
\end{longtable}

\noindent
Las opciones que eligen algo (las glosas, la búsqueda aproximada, la
fuente, los idiomas) lo hacen solo para esa ejecución; los ajustes
conservan lo que tenían.

\subsection{El modo interactivo}

\index{linea de ordenes@línea de órdenes!interactiva}%
En el modo interactivo se busca cada línea que se escribe tras el
indicador \texttt{=>}, salvo si empieza con una de estas:

\begin{longtable}{@{}>{\raggedright}p{0.36\textwidth}p{0.58\textwidth}@{}}
\toprule
Línea & Qué hace\\
\midrule
\endhead
\texttt{\tl E}, \texttt{\tl L} & del español (o el inglés) al latín; del latín\\
\texttt{\tl G} & Logeion o este diccionario\\
\texttt{\tl D [DIC]} & el diccionario de Logeion (solo: una lista para elegir)\\
\texttt{@ARCHIVO} & busca cada línea de ARCHIVO\\
\texttt{\#} & cambia los parámetros (sección~\ref{sec:settings})\\
\texttt{!} & cambia los modos de desarrollador\\
\texttt{\tl H} & la ayuda: cómo buscar y leer los resultados\\
\texttt{\tl A} & acerca de: la versión y los créditos\\
\texttt{?} & la lista de estas líneas\\
\bottomrule
\end{longtable}

\noindent\index{historial de busquedas@historial de búsquedas}%
\key{Arriba} y \key{Abajo} traen las búsquedas anteriores, también las de
sesiones pasadas y las de los otros programas. En los diálogos de
\texttt{\#} y \texttt{!}, \texttt{?} explica cada ajuste. Dos líneas
vacías, o \key{Ctrl+D}, salen. \key{Ctrl+C} no sale: detiene lo que se
está haciendo (una larga lista de resultados, un archivo), y el programa
espera la línea siguiente.

\subsection{El paginador y los colores}

\index{paginador}\index{colores}%
En una terminal, los resultados más largos que la pantalla se abren en el
paginador del sistema, \texttt{less} (o el que nombre \texttt{\$PAGER}),
como con \texttt{man}: \key{RePág} y \key{AvPág} los desplazan,
\texttt{/} busca un texto, y \texttt{q} vuelve al indicador. Los
resultados van en color, con los colores de la terminal, como en la
ventana. Los resultados que se escriben en un archivo o en una tubería
salen sin color, y \texttt{NO\_COLOR=1} quita los colores.

\section{La terminal, a pantalla completa}
\label{sec:tui}

\shot{tui_es}{1}{\texttt{words\_tui.py}, en una terminal.}

\index{terminal a pantalla completa}%
\texttt{words\_tui.py} es la ventana en una terminal
(figura~\ref{fig:tui_es}): una línea de búsqueda arriba, los resultados
debajo y las teclas en la última línea. \key{Arriba} y \key{Abajo} traen
las búsquedas anteriores, y \key{RePág} y \key{AvPág} desplazan los
resultados. En Windows necesita el paquete \texttt{windows-curses}.

\index{terminal a pantalla completa!teclas}%
\begin{longtable}{@{}>{\raggedright}p{0.3\textwidth}p{0.64\textwidth}@{}}
\toprule
Tecla & Acción\\
\midrule
\endhead
\key{F1} & la ayuda, y los créditos\\
\key{F2}, \key{Ctrl+L} & del latín, o hacia el latín\\
\key{F3}, \key{Ctrl+O} & los ajustes (\key{Espacio} cambia uno)\\
\key{F4}, \key{Ctrl+E} & el editor del diccionario\\
\key{F5}, \key{Ctrl+G} & el idioma de los significados (con Logeion, su diccionario)\\
\key{F6} & el idioma de la interfaz\\
\key{F7} & copia todos los resultados\\
\key{F8} & este diccionario o Logeion\\
\key{F9} & borra el historial\\
\key{F10}, \key{Ctrl+Q} & sale\\
\bottomrule
\end{longtable}

\section{La página web}
\label{sec:web}

\shot{web_es}{0.8}{\texttt{words\_web.py}, en el navegador.}

\index{pagina web@página web}%
\texttt{words\_web.py} arranca un pequeño servidor web y abre el navegador
en \url{http://127.0.0.1:8000/}, o en otro puerto si ese está ocupado
(figura~\ref{fig:web_es}). La página tiene el cuadro de búsqueda, la
dirección, \menu{Significados en}, \menu{Búsqueda aproximada},
\menu{Buscar en} y los resultados como en la ventana, con los mismos
enlaces y explicaciones; debajo, las búsquedas anteriores, todos los
ajustes, la ayuda y los créditos. La luna oscurece la página, y el sol la
vuelve a aclarar; hasta que se usa por primera vez, la página sigue al
sistema. Justo antes de la luna, un botón redondo con el código del idioma
de la interfaz (EN, ES o LA) abre la lista de los idiomas.

\index{pagina web@página web!direcciones}%
Cada búsqueda tiene su propia dirección, como
\url{http://127.0.0.1:8000/?q=amicus+est} (\texttt{\&lang=en} para buscar
hacia el latín), así que se puede guardar en marcadores. El servidor
relee el diccionario cuando un editor lo cambia; el diccionario no puede
editarse desde el navegador.

\index{pagina web@página web!opciones}\index{red}%
Solo este equipo puede abrir la página, a menos que \texttt{--host} diga
otra cosa: con \texttt{--host 0.0.0.0} se puede usar desde un teléfono en
la misma red. Como no hay contraseñas, conviene hacerlo solo en una red de
confianza. \texttt{--port} elige el puerto, \texttt{--no-browser} no abre
el navegador, y \key{Ctrl+C} detiene el servidor.

\section{Ajustes}
\label{sec:settings}

\index{ajustes}\index{parametros@parámetros}%
\menu{Diccionario\then Ajustes\ldots} muestra todos los ajustes de la
ventana, con lo que hace cada uno (al poner el puntero sobre un ajuste,
aparece su explicación); un cambio se aplica en el acto. Son los mismos
que los parámetros (\texttt{\#}) y los modos de desarrollador
(\texttt{!}) de la línea de órdenes, \key{F3} en la terminal y los ajustes
de la página web. El diálogo tiene también \menu{Idioma de la interfaz},
\menu{Fuente\ldots}, \menu{Borrar historial} y \menu{Ajustes por defecto}.
Los ajustes conservan su nombre inglés, el mismo del archivo de
configuración.

\index{parametros@parámetros!lo que muestran los resultados}%
\begin{longtable}{@{}>{\raggedright}p{0.36\textwidth}p{0.58\textwidth}@{}}
\toprule
Parámetro & Qué hace, si está activado\\
\midrule
\endhead
\multicolumn{2}{@{}l}{\itshape Lo que muestran los resultados}\\
\texttt{show\_dictionary\_forms} & cada entrada con sus partes principales y su gramática\\
\texttt{show\_only\_meanings} & solo las entradas y sus significados, sin las formas\\
\texttt{use\_glosses} & los análisis en abreviaturas de glosa, no en palabras\\
\texttt{show\_long\_vowels} & las vocales largas marcadas en las formas (\textit{amīc-ōrum})\\
\texttt{show\_codes} & los códigos entre corchetes (\texttt{[XXXBO]})\\
\texttt{explain\_codes} & la época y la frecuencia siempre en palabras, no solo cuando son inusuales\\
\texttt{trim\_output} & omite los resultados menos probables (y, hacia el latín, todos menos los seis mejores)\\
\texttt{trim\_medieval} & con \texttt{trim\_output}, también las formas medievales si se hallaron otras\\
\addlinespace
\multicolumn{2}{@{}l}{\itshape Cómo se buscan las palabras}\\
\texttt{fuzzy\_search} & las grafías cercanas de una palabra que no encuentra nada\\
\texttt{use\_affixes} & una palabra desconocida leída como una conocida con un prefijo o un sufijo\\
\texttt{use\_tricks} & otras grafías, y dos palabras escritas como una\\
\texttt{medieval\_spellings} & con \texttt{use\_tricks}, también las grafías medievales\\
\texttt{read\_compounds} & un participio y una forma de \textit{esse} como tiempo compuesto\\
\texttt{ignore\_unknown\_names} & ningún truco con una palabra desconocida con mayúscula\\
\texttt{near\_entries} & para una palabra desconocida, las entradas cercanas a donde estaría\\
\addlinespace
\multicolumn{2}{@{}l}{\itshape Modos de desarrollador}\\
\texttt{show\_entry\_source} & de qué diccionario es cada entrada, y su número en él\\
\texttt{pearse\_codes} & cada línea de los resultados empieza con un código (para otros programas)\\
\texttt{pause\_in\_screen\_output} & los resultados largos en el paginador, en la línea de órdenes\\
\bottomrule
\end{longtable}

\noindent\index{trim\_output}%
Los poetas usan palabras y formas inusuales, de modo que
\texttt{trim\_output} puede omitir la correcta; con un verso que da
resultados extraños, conviene desactivarlo.

\section{Editar el diccionario}
\label{sec:editing}

\index{editores del diccionario}%
Dos editores añaden, corrigen y borran las entradas del diccionario,
\texttt{data/dictionary.xml}: uno en una ventana,
\texttt{words\_edit\_gui.py}, que también abre \menu{Diccionario\then
Editar diccionario\ldots}, y otro en la terminal, \texttt{words\_edit.py}.
En los dos:

\begin{itemize}
\item una entrada se busca por una forma latina (\textit{amicorum}), una
  raíz o una palabra de su significado;
\item se corrige o se borra, o se hace una nueva, vacía o como la que se
  muestra (\menu{Nueva como esta}); los campos dependen de la categoría
  gramatical, y los códigos se eligen en listas que los explican;
\item la vista previa muestra la entrada como la mostrarán los
  resultados, y cualquier problema que tenga; \menu{Probar una palabra}
  muestra lo que el programa hace con una palabra con el diccionario tal
  como está ahora;
\item también pueden editarse los significados en español, junto a los
  ingleses, tras elegir el idioma con \menu{Significados en} (\key{F7} en
  la terminal).
\end{itemize}

\index{editores del diccionario!guardar}\index{archivos .bak}%
Un cambio entra en el diccionario con \menu{Aplicar cambios} (\key{F3} en
la terminal), y el archivo solo se escribe con \menu{Guardar} (\key{F2}).
Al dejar una entrada con cambios sin aplicar, el editor pregunta qué hacer
con ellos. La primera vez que se guarda en una sesión, los archivos
originales se conservan como archivos \texttt{.bak}.

\index{editores del diccionario!teclas}%
\begin{longtable}{@{}>{\raggedright}p{0.3\textwidth}p{0.64\textwidth}@{}}
\toprule
Tecla (terminal) & Acción\\
\midrule
\endhead
\key{F3}, \key{Ctrl+A} & aplica los cambios (o añade la entrada nueva)\\
\key{F4}, \key{Ctrl+N} & una entrada nueva\\
\key{F5}, \key{Ctrl+D} & una entrada nueva como esta\\
\key{F6}, \key{Ctrl+T} & prueba una palabra\\
\key{F7}, \key{Ctrl+L} & el idioma de los significados editado además del inglés\\
\key{F8}, \key{Ctrl+K} & borra la entrada\\
\key{F2}, \key{Ctrl+S} & guarda el archivo (aplicando antes los cambios)\\
\key{F10}, \key{Ctrl+Q} & sale\\
\key{Tab} & el área siguiente: la búsqueda, la lista, el formulario\\
\key{Ctrl+F} & buscar\\
\bottomrule
\end{longtable}

\noindent\index{editores del diccionario!datos de solo lectura}\index{carpeta de datos}%
En la AppImage y en el RPM los datos no se pueden escribir, así que los
editores están desactivados. Para editar ahí el diccionario, se copia la
carpeta \texttt{data} a otro lugar, y se da la copia como carpeta de
datos: en el archivo de configuración (\texttt{folder} en
\texttt{[data]}), con \texttt{--datadir} o con
\texttt{\$PY\_WORDS\_DATADIR}.

\index{words\_check.py}%
Los archivos también pueden editarse a mano, en cualquier editor de
texto: se copia la entrada de una palabra parecida y se cambia. Luego,
\texttt{python3 words\_check.py} informa, con los números de línea, de lo
que el programa no puede leer y de las entradas que probablemente no hacen
lo que se quería.

\section{Teclas de la ventana}
\label{sec:keys}

\index{teclas}%
\begin{longtable}{@{}>{\raggedright}p{0.4\textwidth}p{0.54\textwidth}@{}}
\toprule
Tecla & Acción\\
\midrule
\endhead
\key{Intro} & buscar\\
\key{Abajo} (en el cuadro de búsqueda) & las búsquedas anteriores\\
\key{Ctrl+F} & buscar en los resultados\\
\key{Intro}, \key{Mayús+Intro} (al buscar en los resultados) & el lugar siguiente y el anterior\\
\key{Esc} (al buscar en los resultados) & cerrar la barra de búsqueda\\
\key{Ctrl++}, \key{Ctrl+-} & texto más grande y más pequeño\\
\key{Ctrl+0} & el tamaño habitual\\
\key{Ctrl+X}, \key{Ctrl+C}, \key{Ctrl+V} & cortar, copiar, pegar\\
\key{F1} & ayuda\\
\key{Ctrl+Q} & salir\\
\key{Esc} (en un diálogo) & cerrarlo\\
\bottomrule
\end{longtable}

\section{Dónde se guarda cada cosa}
\label{sec:files}

\index{archivos de configuracion@archivos de configuración}\index{config.toml}\index{historial de busquedas@historial de búsquedas}\index{cache@caché}%
Los ajustes, las fuentes, el modo oscuro y las opciones de la ventana se
guardan en un solo archivo de configuración, \texttt{config.toml}. Es
texto plano, con un comentario para cada ajuste, y puede editarse a mano;
\texttt{python3 words.py --config} lo crea y dice dónde está. El programa
guarda además las últimas 1000 búsquedas, y una caché del diccionario con
la que arranca en un cuarto de segundo. Todo ello está en las carpetas
del usuario:

\begin{longtable}{@{}llp{0.6\textwidth}@{}}
\toprule
Qué & Sistema & Dónde\\
\midrule
\endhead
Ajustes & Linux & \texttt{\textasciitilde/.config/py-words/config.toml}\\
 & macOS & \texttt{\textasciitilde/Library/Application Support/py-words}\\
 & Windows & \texttt{\%APPDATA\%\textbackslash py-words}\\
\addlinespace
Historial & Linux & \texttt{\textasciitilde/.local/state/py-words/history}\\
 & macOS & \texttt{\textasciitilde/Library/Application Support/py-words}\\
 & Windows & \texttt{\%LOCALAPPDATA\%\textbackslash py-words}\\
\addlinespace
Caché & Linux & \texttt{\textasciitilde/.cache/py-words}\\
 & macOS & \texttt{\textasciitilde/Library/Caches/py-words}\\
 & Windows & \texttt{\%LOCALAPPDATA\%\textbackslash py-words}\\
\bottomrule
\end{longtable}

\noindent\index{variables de entorno}%
\texttt{\$PY\_WORDS\_CONFIG}, \texttt{\$PY\_WORDS\_HISTORY} y
\texttt{\$PY\_WORDS\_CACHE} dan otros lugares; con
\texttt{\$PY\_WORDS\_HISTORY} vacía no se guarda historial, y con
\texttt{\$PY\_WORDS\_CACHE} vacía no hay caché. La caché siempre puede
borrarse. Los archivos de datos se buscan primero con
\texttt{--datadir}, luego en \texttt{\$PY\_WORDS\_DATADIR}, luego en la
carpeta que nombre el archivo de configuración y por último en la carpeta
\texttt{data} junto al programa.

\section{Cuando algo falla}

\index{problemas}%
\begin{description}
\item[El programa dice que Python es demasiado antiguo] instale Python
  3.13 o posterior, o use la AppImage.
\item[La ventana no se abre: no module named tkinter] instale Tk
  (sección~\ref{sec:install}).
\item[Una palabra no se encuentra] compruebe \menu{Búsqueda aproximada}
  y los ajustes \texttt{use\_tricks} y \texttt{medieval\_spellings}; una
  palabra desconocida con mayúscula se toma como un nombre, salvo que
  \texttt{ignore\_unknown\_names} esté desactivado.
\item[Un verso da resultados extraños] desactive \texttt{trim\_output},
  que omite las formas inusuales.
\item[Logeion no responde] necesita internet; \menu{Buscar en\then Local}
  vuelve a este diccionario.
\item[\menu{Editar diccionario\ldots} está en gris] los datos no se pueden
  escribir (AppImage, RPM), o está elegido Logeion
  (sección~\ref{sec:editing}).
\item[El texto de la ventana es muy pequeño o muy grande] \key{Ctrl++} y
  \key{Ctrl+-}, o \menu{Ver\then Fuente\ldots}.
\item[Un ajuste quedó mal] borre \texttt{config.toml}
  (sección~\ref{sec:files}); el programa arranca con sus valores por
  defecto.
\end{description}

\noindent
Los problemas y los deseos pueden escribirse como \textit{issues} en
GitHub o en Codeberg (las direcciones están en la primera página), o
enviarse a
\href{mailto:nicolasvaughan@proton.me}{\texttt{nicolasvaughan@proton.me}}.

\printindex

\end{document}

% ===============================================
%%% Local Variables:
%%% mode: latex
%%% lsp-ltex-plus-language: "es-CO"
%%% TeX-master: t
%%% End:
